\documentclass[11pt,a4paper]{article}

% --- Encoding and fonts ---
\usepackage[utf8]{inputenc}
\usepackage[T1]{fontenc}

% --- Mathematics ---
\usepackage{amsmath,amssymb,amsthm}
\usepackage{mathtools}

% --- Page layout ---
\usepackage[margin=1in,headheight=14pt]{geometry}

% --- Hyperlinks ---
\usepackage[colorlinks=true,linkcolor=red,citecolor=blue,urlcolor=blue]{hyperref}

% --- Lists ---
\usepackage{enumitem}

% --- Colors and boxes ---
\usepackage{xcolor}
\usepackage[theorems,skins,breakable]{tcolorbox}

% --- Header/footer ---
\usepackage{fancyhdr}
\pagestyle{fancy}
\fancyhf{}
\fancyhead[L]{\textsc{Lesson 9: Convex Optimization and Duality}}
\fancyhead[R]{\thepage}

% --- Tables ---
\usepackage{booktabs}

% ============================================================
% Theorem environments
% ============================================================
\theoremstyle{definition}
\newtheorem{definition}{Definition}[section]
\newtheorem{example}{Example}[section]
\newtheorem{exercise}{Exercise}[section]

\theoremstyle{plain}
\newtheorem{theorem}{Theorem}[section]
\newtheorem{lemma}[theorem]{Lemma}
\newtheorem{proposition}[theorem]{Proposition}
\newtheorem{corollary}[theorem]{Corollary}

\theoremstyle{remark}
\newtheorem{remark}{Remark}[section]

% ============================================================
% Colored boxes
% ============================================================
\newtcolorbox{keyresult}[1][]{
  colback=green!5!white,
  colframe=green!60!black,
  fonttitle=\bfseries,
  title={Key Result},
  breakable,
  #1
}

\newtcolorbox{rlconnection}[1][]{
  colback=blue!5!white,
  colframe=blue!60!black,
  fonttitle=\bfseries,
  title={RL Connection},
  breakable,
  #1
}

\newtcolorbox{warning}[1][]{
  colback=red!5!white,
  colframe=red!60!black,
  fonttitle=\bfseries,
  title={Warning},
  breakable,
  #1
}

% ============================================================
% Custom commands
% ============================================================
\newcommand{\R}{\mathbb{R}}
\newcommand{\N}{\mathbb{N}}
\newcommand{\Q}{\mathbb{Q}}
\newcommand{\Z}{\mathbb{Z}}
\newcommand{\C}{\mathbb{C}}
\newcommand{\Prob}{\mathbb{P}}
\newcommand{\E}{\mathbb{E}}
\newcommand{\indicator}{\mathbf{1}}

\DeclareMathOperator{\dom}{dom}
\DeclareMathOperator{\epi}{epi}
\DeclareMathOperator{\conv}{conv}
\DeclareMathOperator{\prox}{prox}
\DeclareMathOperator*{\argmin}{arg\,min}
\DeclareMathOperator*{\argmax}{arg\,max}
\DeclareMathOperator{\tr}{tr}
\DeclareMathOperator{\diag}{diag}
\DeclareMathOperator{\interior}{int}

% ============================================================
\title{Lesson 26: Lagrangian Duality and KKT Conditions}
\author{%
  Session 9 of 102 \quad$\mid$\quad Phase I: Mathematical Foundations \quad$\mid$\quad 8h \\[4pt]
  \textit{Primary texts: Boyd \& Vandenberghe Ch.\,4--5, 9--11; Nesterov Ch.\,1--4}
}
\date{}

\begin{document}
\maketitle
\tableofcontents
\newpage

\setcounter{section}{2}

\section{Lagrangian Duality}\label{sec:duality}

\subsection{The Lagrangian and Dual Function}

\begin{definition}[Lagrangian]\label{def:lagrangian}
Given the convex program~\eqref{eq:std-form}, the \emph{Lagrangian}
$L : \R^n \times \R^m \times \R^p \to \R$ is
\begin{equation}\label{eq:lagrangian}
  L(x, \lambda, \nu) = f_0(x) + \sum_{i=1}^m \lambda_i f_i(x)
    + \sum_{j=1}^p \nu_j (a_j^\top x - b_j),
\end{equation}
where $\lambda \in \R^m_+$ are the \emph{dual variables} (or \emph{Lagrange
multipliers}) for the inequality constraints and $\nu \in \R^p$ are the dual
variables for the equality constraints.
\end{definition}

\begin{definition}[Lagrange dual function]\label{def:dual-fn}
The \emph{Lagrange dual function} $g : \R^m \times \R^p \to
\R \cup \{-\infty\}$ is
\begin{equation}\label{eq:dual-fn}
  g(\lambda, \nu) = \inf_{x \in \R^n} L(x, \lambda, \nu).
\end{equation}
\end{definition}

\begin{proposition}[Dual function is concave]\label{prop:dual-concave}
For any optimization problem (not necessarily convex), the dual function $g$ is
concave, even when the primal problem is not convex.
\end{proposition}

\begin{proof}
For each fixed $x$, the map $(\lambda, \nu) \mapsto L(x, \lambda, \nu)$ is
affine (hence concave) in $(\lambda, \nu)$. The pointwise infimum over $x$ of
a family of concave functions is concave.
\end{proof}

\subsection{Weak Duality}

\begin{definition}[Dual problem]\label{def:dual-problem}
The \emph{Lagrange dual problem} is
\begin{equation}\label{eq:dual-problem}
  \max_{\lambda \ge 0,\, \nu} \; g(\lambda, \nu).
\end{equation}
Its optimal value is denoted $d^*$.
\end{definition}

\begin{theorem}[Weak duality]\label{thm:weak-duality}
For any optimization problem (not necessarily convex),
\begin{equation}\label{eq:weak-duality}
  d^* \le p^*.
\end{equation}
The quantity $p^* - d^*$ is called the \emph{duality gap}.
\end{theorem}

\begin{proof}
Let $\tilde{x}$ be any feasible point for the primal problem, so
$f_i(\tilde{x}) \le 0$ and $a_j^\top \tilde{x} = b_j$. For any
$\lambda \ge 0$ and any $\nu$,
\[
  L(\tilde{x}, \lambda, \nu)
  = f_0(\tilde{x}) + \sum_{i=1}^m \underbrace{\lambda_i}_{\ge 0}
    \underbrace{f_i(\tilde{x})}_{\le 0}
  + \sum_{j=1}^p \nu_j \underbrace{(a_j^\top \tilde{x} - b_j)}_{= 0}
  \le f_0(\tilde{x}).
\]
Taking the infimum over $x$ gives $g(\lambda, \nu) \le L(\tilde{x}, \lambda, \nu) \le f_0(\tilde{x})$.
Since this holds for every feasible $\tilde{x}$, we get $g(\lambda, \nu) \le p^*$.
Maximizing the left side over $(\lambda, \nu)$ with $\lambda \ge 0$ yields
$d^* \le p^*$.
\end{proof}

\subsection{Strong Duality and Slater's Condition}

\begin{definition}[Strong duality]\label{def:strong-duality}
\emph{Strong duality} holds when $d^* = p^*$, i.e., the duality gap is zero.
\end{definition}

\begin{definition}[Slater's condition]\label{def:slater}
A convex program~\eqref{eq:std-form} satisfies \emph{Slater's condition}
(or \emph{Slater's constraint qualification}) if there exists a point
$\hat{x} \in \text{relint}(\dom f_0)$ such that
\[
  f_i(\hat{x}) < 0, \quad i = 1,\dots,m, \qquad
  a_j^\top \hat{x} = b_j, \quad j = 1,\dots,p.
\]
Such $\hat{x}$ is called a \emph{strictly feasible point}.
\end{definition}

\begin{theorem}[Slater implies strong duality]\label{thm:strong-duality}
If the convex program~\eqref{eq:std-form} is feasible and Slater's condition
holds, then strong duality holds: $d^* = p^*$. Moreover, if $p^*$ is finite,
the dual optimum is attained.
\end{theorem}

\begin{proof}[Proof sketch]
We use a supporting-hyperplane argument on the set
\[
  \mathcal{G} = \{(u, v, t) \in \R^m \times \R^p \times \R :
    \exists\, x \text{ with } f_i(x) \le u_i,\; a_j^\top x - b_j = v_j,\;
    f_0(x) \le t\}.
\]
This set is convex (the $f_i$ are convex and $a_j^\top x - b_j$ is affine).
The point $(0, 0, p^*)$ lies on the boundary of $\mathcal{G}$ (it is achieved
in the limit by a minimizing sequence).

We claim $(0, 0, p^*-\epsilon) \notin \mathcal{G}$ for any $\epsilon > 0$, by
definition of $p^*$. Hence there is a supporting hyperplane at $(0, 0, p^*)$:
there exist $(\lambda, \nu, \mu) \ne 0$ with $\mu \ge 0$ such that
\[
  \lambda^\top u + \nu^\top v + \mu t \ge \mu p^*
\]
for all $(u, v, t) \in \mathcal{G}$, with $\lambda \ge 0$.

\emph{Claim:} $\mu > 0$. Suppose $\mu = 0$. Then
$\lambda^\top u + \nu^\top v \ge 0$ for all $(u,v,t) \in \mathcal{G}$.
Evaluating at $x = \hat{x}$ (the Slater point) with $u_i = f_i(\hat{x}) < 0$,
$v_j = 0$, gives $\sum_i \lambda_i f_i(\hat{x}) \ge 0$. Since $\lambda \ge 0$
and $f_i(\hat{x}) < 0$, we get $\lambda = 0$. Then $\nu^\top v \ge 0$ for all
achievable $v$, which forces $\nu = 0$ (since $v$ can be any value via
adjusting $x$). This contradicts $(\lambda, \nu, \mu) \ne 0$.

With $\mu > 0$, divide through by $\mu$ and set $\bar{\lambda} = \lambda/\mu$,
$\bar{\nu} = \nu/\mu$. For any feasible $x$ (where $u_i = f_i(x)$, $v_j = 0$,
$t = f_0(x)$):
\[
  \sum_i \bar{\lambda}_i f_i(x) + f_0(x) \ge p^*.
\]
Taking the infimum over $x$ gives
$g(\bar{\lambda}, \bar{\nu}) \ge p^*$. Combined with weak duality $g \le p^*$,
we conclude $d^* \ge g(\bar{\lambda}, \bar{\nu}) = p^*$, hence $d^* = p^*$.
The dual optimum is attained at $(\bar{\lambda}, \bar{\nu})$.
\end{proof}

\begin{rlconnection}
Strong duality is the theoretical foundation for the Lagrangian approach to
constrained MDPs (CMDPs). In a CMDP, we augment the reward objective with
penalty terms for constraint violations. Slater's condition guarantees there is
no duality gap, so solving the dual (which is unconstrained) is equivalent to
solving the primal. This justifies primal-dual policy optimization algorithms
such as Lagrangian PPO for safe RL.
\end{rlconnection}

%% ============================================================
%% SECTION 4: KKT Conditions
%% ============================================================

\section{Karush--Kuhn--Tucker Conditions}\label{sec:kkt}

\begin{definition}[KKT conditions]\label{def:kkt}
For the convex program~\eqref{eq:std-form}, the \emph{Karush--Kuhn--Tucker
(KKT) conditions} for the primal-dual triple $(x^*, \lambda^*, \nu^*)$ are:
\begin{enumerate}[label=(KKT\arabic*)]
  \item \textbf{Primal feasibility:}
    $f_i(x^*) \le 0$ for $i = 1,\dots,m$; \;
    $a_j^\top x^* = b_j$ for $j = 1,\dots,p$.
  \item \textbf{Dual feasibility:}
    $\lambda^*_i \ge 0$ for $i = 1,\dots,m$.
  \item \textbf{Complementary slackness:}
    $\lambda^*_i f_i(x^*) = 0$ for $i = 1,\dots,m$.
  \item \textbf{Stationarity:}
    $\nabla f_0(x^*) + \sum_{i=1}^m \lambda^*_i \nabla f_i(x^*)
     + \sum_{j=1}^p \nu^*_j a_j = 0$.
\end{enumerate}
\end{definition}

\begin{theorem}[KKT necessary conditions under strong duality]\label{thm:kkt-necessary}
If strong duality holds and $x^*$ is primal optimal, $(\lambda^*, \nu^*)$ is
dual optimal, and the $f_i$ are differentiable at $x^*$, then the KKT
conditions hold at $(x^*, \lambda^*, \nu^*)$.
\end{theorem}

\begin{proof}
Since strong duality holds, $f_0(x^*) = p^* = d^* = g(\lambda^*, \nu^*)$.
By definition,
\[
  g(\lambda^*, \nu^*) = \inf_x L(x, \lambda^*, \nu^*) \le L(x^*, \lambda^*, \nu^*).
\]
We expand:
\[
  L(x^*, \lambda^*, \nu^*)
  = f_0(x^*) + \sum_i \lambda^*_i f_i(x^*) + \sum_j \nu^*_j(a_j^\top x^* - b_j).
\]
By primal feasibility, $a_j^\top x^* - b_j = 0$ and $f_i(x^*) \le 0$. Since
$\lambda^* \ge 0$, we have $\sum_i \lambda^*_i f_i(x^*) \le 0$, so
$L(x^*, \lambda^*, \nu^*) \le f_0(x^*)$.

Combining: $f_0(x^*) = g(\lambda^*, \nu^*) \le L(x^*, \lambda^*, \nu^*) \le f_0(x^*)$.
Hence all inequalities are equalities. In particular,
$\sum_i \lambda^*_i f_i(x^*) = 0$. Since each term is non-positive
($\lambda^*_i \ge 0$, $f_i(x^*) \le 0$), each term is zero:
$\lambda^*_i f_i(x^*) = 0$. This gives \textbf{complementary slackness}.

Furthermore, $g(\lambda^*, \nu^*) = L(x^*, \lambda^*, \nu^*)$, meaning $x^*$
minimizes $L(\cdot, \lambda^*, \nu^*)$ over $\R^n$. Since this is an
unconstrained minimum of a differentiable function, the gradient vanishes:
\[
  \nabla_x L(x^*, \lambda^*, \nu^*)
  = \nabla f_0(x^*) + \sum_i \lambda^*_i \nabla f_i(x^*) + \sum_j \nu^*_j a_j
  = 0.
\]
This is \textbf{stationarity}. Primal and dual feasibility hold by assumption.
\end{proof}

\begin{theorem}[KKT sufficient conditions for convex problems]\label{thm:kkt-sufficient}
For the convex program~\eqref{eq:std-form} with differentiable $f_0, \dots, f_m$,
if a triple $(x^*, \lambda^*, \nu^*)$ satisfies the KKT conditions, then $x^*$
is primal optimal and $(\lambda^*, \nu^*)$ is dual optimal, with zero duality
gap.
\end{theorem}

\begin{proof}
By stationarity, $x^*$ minimizes $L(\cdot, \lambda^*, \nu^*)$ (since
$L$ is convex in $x$ and the gradient vanishes at $x^*$). Therefore,
\[
  g(\lambda^*, \nu^*) = \inf_x L(x, \lambda^*, \nu^*) = L(x^*, \lambda^*, \nu^*).
\]
By complementary slackness and primal feasibility,
\[
  L(x^*, \lambda^*, \nu^*) = f_0(x^*) + \underbrace{\sum_i \lambda^*_i f_i(x^*)}_{=\,0}
  + \underbrace{\sum_j \nu^*_j (a_j^\top x^* - b_j)}_{=\,0} = f_0(x^*).
\]
Hence $g(\lambda^*, \nu^*) = f_0(x^*)$. By weak duality,
$d^* \ge g(\lambda^*, \nu^*) = f_0(x^*) \ge p^* \ge d^*$, so all are equal.
Thus $x^*$ is primal optimal and $(\lambda^*, \nu^*)$ is dual optimal.
\end{proof}

\begin{keyresult}[title={KKT Conditions for Convex Problems}]
For convex programs with differentiable objectives and constraints, the KKT
conditions are \textbf{necessary} (assuming a constraint qualification such as
Slater's condition, which implies strong duality) and \textbf{sufficient} for
optimality. This makes them the central tool for characterizing solutions.
\end{keyresult}

\begin{rlconnection}
The KKT conditions characterize solutions of the policy optimization
subproblem in TRPO and PPO. When we constrain the policy update to a KL
trust region, the KKT multiplier on the KL constraint determines the
effective step size. Complementary slackness tells us: either the constraint
is active (the update reaches the trust-region boundary) or the multiplier is
zero (the unconstrained optimum lies inside the trust region).
\end{rlconnection}

%% ============================================================
%% SECTION 5: Gradient Descent
%% ============================================================


\end{document}